\scanpage{111}
\mspage{\handwritten{86}}

Utilizing (14) and assuming $n$ to be sufficiently
large, we find from (21) and (22):

\[
\handwritten{\dim H^1(V,O_V[nE'])=\dim H^1(V,O_V(-\mathscr E'_\nu)[nE']),}\tag*{\handwritten{(23)}}
\]

We can assume that $m'\geq n(\handwritten{\mathscr E'}_\nu)$ and so great that,
at $n=m'$, (23) is fulfilled. Then we have an exact
sequence

\diagram{\handwritten{\begin{tikzcd}[ampersand replacement=\&]
H^0(V,O_V[mE'])\arrow[r,"{r^*}"] \& H^0(\mathscr E'_\nu,O_{\mathscr E'_\nu}[m'E'])\arrow[r] \& 0
\end{tikzcd}}}

Utilizing the properties of the homomorphism
$r^*$ (See p. 208) we obtain hence that the mapping of the
variety $\handwritten{\mathscr E'}_\nu$ into a projective space, corresponding to
the group $H^{\unsure{0|O}}\handwritten{\big(}(\handwritten{\mathscr E'}_\nu),O_{\handwritten{\mathscr E'}_\nu}[m'E'])_{\handwritten{\mathscr E'}_\nu}\handwritten{\big)}$, coincides with the
restriction on $\handwritten{\mathscr E'}_\nu$ of the mapping $f_{m'}$. In particular,
we obtain that $f_{m'}$ maps $\handwritten{\mathscr E'}_\nu$ in a one-to-one manner
(that is, different points are transfered into different
points).

Let $q'$ be so great that $r=2^{q'}m-m'>0$. Let
us show that $f_{2^{q'}m}$ maps $\handwritten{\mathscr E'}_\nu$ in a one-to-one manner.
Let $P_1$ and $P_2$ be any two points on $\handwritten{\mathscr E'}_\nu$. It is
sufficient for us to show that there exists a cross-section of $H^0(V,O_V[2^{q'}mE'])$, which is equal to zero
in $P_1$, but not in $P_2$. It follows from above that there

% archive sheet number: 110
